\documentclass[11pt,a4paper]{article}
\usepackage{curso}
\begin{document}
\portada{Lección}{3.6}{Relaciones de incertidumbre termodinámica}

\noindent\textbf{Unidad 3.}\quad\textbf{Objetivos:} enunciar la TUR y la cota cinética; compararlas; aplicarlas al tiempo de cómputo.
\textbf{Prerrequisitos:} 3.2 (producción de entropía), 2.8 (procesos de saltos y Gillespie).

\section{Corrientes y precisión}
Una \emph{corriente} $J$ es un observable antisimétrico bajo inversión temporal que crece con el tiempo: pasos netos de un motor, tics de
un reloj, tokens procesados por un circuito browniano. Su precisión se mide con la varianza relativa
$\epsilon^2 = \mathrm{Var}(J)/\langle J\rangle^2$.

\begin{intuicion}
Un contador que salta adelante con tasa $k_+$ y atrás con tasa $k_-$ disipa $\ln(k_+/k_-)$ (en unidades de $\kB$) por cada paso neto.
Si $k_+\approx k_-$ casi no disipa, pero avanza de forma errática; si $k_-\ll k_+$ avanza con regularidad, pero cada paso cuesta mucha entropía.
\end{intuicion}

\section{La relación de incertidumbre termodinámica}
\begin{resultado}[ (TUR)]
Para un proceso de Markov en estado estacionario, toda corriente cumple
\begin{equation}
\frac{\mathrm{Var}(J)}{\langle J\rangle^2}\ \ge\ \frac{2\kB}{\Sigma},
\end{equation}
donde $\Sigma$ es la entropía total producida en el intervalo de observación.
\end{resultado}
\noindent Para el contador sesgado en tiempo $t$: $\langle J\rangle=(k_+-k_-)t$, $\mathrm{Var}(J)=(k_++k_-)t$ y $\Sigma=(k_+-k_-)t\ln(k_+/k_-)$, de modo que
\begin{equation}
\epsilon^2\Sigma=\frac{(k_++k_-)\ln(k_+/k_-)}{k_+-k_-}\ \ge\ 2,
\end{equation}
con igualdad en el límite $k_+\to k_-$ (desarrollo: $\ln(1+x)\approx x-x^2/2$). Reducir el error relativo a la mitad exige cuadruplicar $\Sigma$.

\section{La cota cinética}
\begin{resultado}[ (cota cinética)]
Con la actividad dinámica $\mathcal A$ (número medio total de saltos, adelante y atrás), $\ \mathrm{Var}(J)/\langle J\rangle^2\ \ge\ 1/\mathcal A$.
\end{resultado}
Para el contador, $\mathcal A=(k_++k_-)t$ y $\epsilon^2\mathcal A=(k_++k_-)^2/(k_+-k_-)^2\ge1$, ajustada cuando $k_-\to0$. La TUR es ajustada cerca
del equilibrio y la cinética lejos de él; existen cotas mixtas que combinan $\Sigma$ y $\mathcal A$.

\section{Tiempo de cómputo}
El tiempo $T_N$ hasta que el contador llega a $N$ pasos netos es un tiempo de primer paso. Para $N$ grande,
$\langle T_N\rangle\approx N/(k_+-k_-)$ y $\epsilon_T^2\approx(k_++k_-)/[N(k_+-k_-)]$; las mismas cotas valen con $\Sigma\approx N\ln(k_+/k_-)$ y
$\mathcal A\approx(k_++k_-)\langle T_N\rangle$ \cite{utsumi}.
\begin{ejemplo}
$N=100$, $k_-=1$, simulación de Gillespie \evS:
\begin{center}\begin{tabular}{lccc}\toprule
$k_+$ & $\epsilon_T^2$ & $2/\Sigma$ & $1/\mathcal A$\\\midrule
1,2 & 0,110 & 0,110 & 0,0009\\
3 & 0,020 & 0,018 & 0,005\\
20 & 0,011 & 0,0067 & 0,009\\\bottomrule\end{tabular}\end{center}
Con $k_+=1{,}2$ la TUR está saturada; con $k_+=20$ la cota cinética es la más ajustada. Ambas se cumplen siempre.
\end{ejemplo}

\begin{aplicacion}
Circuitos brownianos de tokens (lección 5.5, reto L7): el tiempo de cómputo de una compuerta obedece una TUR \cite{utsumi}. Inferencia
termodinámica: medir $\langle J\rangle$ y $\mathrm{Var}(J)$ da la cota inferior $\Sigma\ge2\langle J\rangle^2/\mathrm{Var}(J)$ sin conocer el modelo.
\end{aplicacion}

\begin{error}
\begin{itemize}[nosep]
\item Aplicar la TUR fuera de Markov estacionario sin la versión adecuada. \quad Usar la \emph{tasa} de entropía en lugar de $\Sigma$ total.
\item Olvidar el factor 2 (con $\kB=1$). \quad Suponer que la cota siempre está saturada.
\end{itemize}
\end{error}

\begin{thebibliography}{9}\small
\bibitem[26]{utsumi} Y. Utsumi et al., Computation time and thermodynamic uncertainty relation of Brownian circuits, 2022.
\bibitem[2]{pp} L. Peliti y S. Pigolotti, \emph{Stochastic Thermodynamics: An Introduction}, Princeton, 2021.
\bibitem[11]{seifert} U. Seifert, Stochastic thermodynamics, fluctuation theorems and molecular machines, Rep. Prog. Phys. 75 (2012).
\end{thebibliography}
\end{document}
